% =============================================================================
% Section 3: Scenario Design
% =============================================================================
\section{Scenario Design}
\label{sec:scenarios}

To disentangle the relative contributions of climate forcing and pathogen
evolution to the long-term trajectory of \textit{Pycnopodia helianthoides}
populations, we designed a factorial scenario ensemble that independently
varied two axes: the magnitude of ocean warming and the availability of
evolutionary mechanisms within the pathogen module.  All scenarios were
executed with the validated production model W154 (coast-wide
RMSE\,=\,0.504; Section~\ref{sec:calibration}).

\subsection{Factorial Structure}

The four forecast scenarios (Table~\ref{tab:scenarios}) arise from the
cross of two climate pathways and two evolutionary configurations:

\begin{table}[htbp]
\centering
\caption{Factorial scenario design.  ``Evo ON'' enables both pathogen
thermal-niche adaptation and local virulence evolution; ``Evo OFF''
freezes all pathogen traits at their calibrated 2012 values.  F03
provides an intermediate case in which the thermal niche may shift but
virulence remains static.}
\label{tab:scenarios}
\smallskip
\begin{tabular}{llccc}
\toprule
\textbf{ID} & \textbf{Climate} & \textbf{Thermal adapt.} &
  \textbf{Virulence evol.} & \textbf{Rationale} \\
\midrule
F01 & SSP2-4.5 & ON  & ON  & Baseline: moderate warming, full evolution \\
F02 & SSP2-4.5 & OFF & OFF & Isolate total evolutionary effect \\
F03 & SSP2-4.5 & ON  & OFF & Isolate virulence evolution contribution \\
F04 & SSP5-8.5 & ON  & ON  & Probe sensitivity to high-end warming \\
\bottomrule
\end{tabular}
\end{table}

\noindent
The baseline scenario F01 represents the ``most likely'' coupled
trajectory: a middle-of-the-road emissions pathway (SSP2-4.5) with all
eco-evolutionary feedbacks active.  Comparing F01 to F02 quantifies the
aggregate effect of pathogen evolution, while comparing F01 to F03
isolates the marginal contribution of virulence evolution alone (with
thermal adaptation still permitted).  Finally, F04 retains all
evolutionary mechanisms but substitutes the high-emissions SSP5-8.5
climate trajectory, enabling direct assessment of the warming signal
independent of evolutionary configuration.

This design permits three clean contrasts:
\begin{enumerate}
  \item \textbf{Total pathogen evolution effect:} F02\,$-$\,F01
        (both climate and evolutionary axes vary);
  \item \textbf{Virulence evolution alone:} F01\,$-$\,F03
        (thermal adaptation held constant, virulence toggled);
  \item \textbf{Climate forcing:} F04\,$-$\,F01
        (evolution held constant, SSP pathway varied).
\end{enumerate}

\subsection{Common Model Settings}

All four scenarios share the following configuration:
\begin{itemize}
  \item \textbf{Spatial domain:} 896 coastal grid cells spanning the
        full historical range of \textit{P.~helianthoides} from the
        western Aleutian Islands (52.5\textdegree N) to central Baja
        California (28.5\textdegree N).
  \item \textbf{Carrying capacity:} $K = 5{,}000$ individuals per cell,
        reflecting the maximum local density supported by prey
        availability and habitat suitability.
  \item \textbf{Model version:} W154 production parameterization, as
        selected through the calibration procedure described in
        Section~\ref{sec:calibration}.
  \item \textbf{Temporal extent:} 2012--2050, with daily time steps.
  \item \textbf{Replication:} Each scenario was run with three
        independent random seeds to characterize stochastic variability.
        All results reported below are three-seed means unless otherwise
        noted.
\end{itemize}

\subsection{Sea Surface Temperature Forcing}
\label{sec:sst_forcing}

The thermal forcing protocol distinguishes two periods.  During the
\textit{observation window} (2012--2025), all scenarios use
daily sea surface temperature (SST) fields derived from the NOAA
Optimum Interpolation SST v2.1 dataset \citep{reynolds2007,
huang2021}, bilinearly interpolated to each of the 896 model grid
cells.  This ensures that the simulated SSWD epizootic and its
immediate aftermath are driven by observed oceanographic conditions,
including the anomalous 2013--2016 marine heatwave.

For the \textit{projection window} (2026--2050), SSTs are generated
via a delta method applied to an ensemble of CMIP6 general circulation
models.  Eleven reference coastal sites, distributed at regular
latitudinal intervals from Alaska to Baja California, were selected to
span the major oceanographic provinces within the model domain.  At
each reference site, monthly climatological differences between
historical and future SSP-forced CMIP6 runs were computed and smoothed
with a 12-month running mean to remove interannual noise while
preserving the secular warming trend.  These monthly deltas were then
spatially interpolated (inverse-distance weighting from the nearest
three reference sites) and added to the 2012--2025 OISST climatology
at each of the 896 model cells.  A stochastic daily perturbation term,
calibrated to match the observed interannual SST variance at each
site, was superimposed to maintain realistic thermal variability.

Under SSP2-4.5, this procedure yields coast-wide mean warming of
approximately $+$0.8\,\textdegree C by 2050 relative to the
2012--2020 baseline, with stronger warming in the Gulf of Alaska and
weaker warming along the upwelling-dominated California Current.
Under SSP5-8.5 (scenario F04), warming roughly doubles to
$+$1.5\,\textdegree C coast-wide, with pronounced amplification at
high latitudes.

\subsection{Initialization and Spin-Up}

All scenarios are initialized identically at 1 January 2012 using the
pre-epizootic abundance estimates described in Section~\ref{sec:abm_framework}.
Host resistance allele frequencies are set to $r_0 = 0.150$ uniformly
across all populations, reflecting the calibrated standing genetic
variation prior to the onset of selection by SSWD.  Pathogen traits
($T_\mathrm{vbnc}$, $v_\mathrm{local}$) are likewise initialized at
their calibrated 2012 values.  No burn-in period is applied; the
epizootic itself (2013--2015) serves as the dynamical transition from
pre-disease equilibrium to the post-outbreak regime that is the focus
of the forecast analysis.
